\begin{table}[tp]
\centering
\caption{Cyclic gradients for \href{https://github.com/k3nfalt/longtails/blob/main/LongTails/Setups/InterpolatedFiniteSum.lean}{InterpolatedFiniteSum}.}\label{tab:InterpolatedFiniteSum}
\parbox{\linewidth}{\footnotesize On a real Hilbert space, $m\ge1$ convex differentiable components have $L$-Lipschitz gradients, $L>0$, and a common minimizer $x_\star$ with $\|x_0-x_\star\|\le R$, $R\ge0$. Exact component gradients are available. The setup fixes cyclic gradient updates with step $1/L$ and the last complete-pass iterate; the target is $F(x)-F(x_\star)\le\varepsilon$, $\varepsilon>0$, almost surely. Updates use $m,L,x_0$; the stopping budget also uses $R,\varepsilon$.}\par\smallskip
\begin{threeparttable}
\footnotesize
\setlength{\tabcolsep}{4pt}\renewcommand{\arraystretch}{1.25}
\begin{tabularx}{\linewidth}{@{}>{\raggedright\arraybackslash}p{0.20\linewidth}>{\raggedright\arraybackslash}p{0.32\linewidth}>{\raggedright\arraybackslash}p{0.25\linewidth}>{\raggedright\arraybackslash}X@{}}
\toprule
\textbf{Result} & \textbf{Oracle complexity} & \textbf{Novelty} & \textbf{Limitations} \\
\midrule
\textbf{(1)} \href{https://github.com/k3nfalt/longtails/blob/main/LongTails/Theorems/CyclicGradientInterpolation.lean}{Cyclic\allowbreak{}Gradient\allowbreak{}Interpolation} \textbf{(Ours)} \href{https://github.com/k3nfalt/longtails/blob/main/scripts/reviews/CyclicGradientInterpolation.md}{(review)} & $\mathcal{O}\!\left(m+\frac{m^4LR^2}{\varepsilon}\right)$ & --- & --- \\
\midrule[1.2pt]
FISTA on $F$ \citep[Thm.~4.4]{ifsBeckTeboulle2009} & $\mathcal{O}\!\left(m+m\sqrt{\frac{LR^2}{\varepsilon}}\right)$\par Finite-dimensional restriction; $m$ component calls per step. & \textbf{(1)}: no established improvement in oracle complexity. & \textbf{(1)}: $m^4LR^2/\varepsilon$ instead of $m\sqrt{LR^2/\varepsilon}$; requires interpolation, unnecessary for this baseline. \\
\midrule
Cyclic IGD \citep[Thm.~2.3]{ifsCaiDiakonikolas2024v2} & $\mathcal{O}\!\left(m+\frac{mLR^2\sqrt{\ell}}{\varepsilon}\right)$\par Finite-dimensional, interpolating restriction; horizon-tuned step. & i) \textbf{(1)}: Hilbert-space scope.\par ii) \textbf{(1)}: removes $\sqrt{\ell}$, but costs $m^3$ more; rate gain only for $\ell\gg m^6$. & \textbf{(1)}: worse leading term when $\ell\ll m^6$; requires interpolation, unnecessary for the general baseline. \\
\midrule
Cyclic projections, derived bound \citep[Prop.~4.13]{ifsEvronEtAl2023} & $\mathcal{O}\!\left(m+\frac{m^4L^2R^4}{\varepsilon^2}\right)$\par $m\ge3$, finite dimension; $f_i=\tfrac L2\operatorname{dist}(\cdot,C_i)^2$, closed convex $C_i$ with nonempty intersection. & i) \textbf{(1)}: general smooth interpolating losses.\par ii) \textbf{(1)}: same cyclic method, $\varepsilon^{-1}$ instead of $\varepsilon^{-2}$ on this subclass when $\varepsilon<LR^2$. & \textbf{(1)}: average loss, rather than maximum squared distance; projection specialization is not separately formalized. \\
\bottomrule
\end{tabularx}
\begin{tablenotes}[flushleft]\footnotesize
\item[] \textbf{Abbr:} $F=m^{-1}\sum_{i=1}^m f_i$; $m$: component count; $L$: component smoothness; $R$: initial-distance bound; $\varepsilon$: objective accuracy; $\ell=1+\log(1+LR^2/\varepsilon)$. Costs count component-gradient calls; on squared-distance losses, one call equals one projection. IGD: incremental gradient descent.
\end{tablenotes}
\end{threeparttable}
\end{table}
\clearpage
